% Libro 4 · Mundo 5 · Lenguaje figurado (app: pu-figurado)
\mundo{Lenguaje figurado}{pu-figurado}{Cuando un texto dice «el tiempo es oro», nadie piensa en joyas. El lenguaje
figurado dice una cosa para significar otra, y es tan común en los textos literarios como en los periodísticos. El examen
pide reconocer las figuras y, sobre todo, entender qué significan en su contexto, incluida la más traicionera: la ironía.}{%
\item \textbf{5.1} Figuras literarias
\item \textbf{5.2} Ironía}

\leccion{pu-lf-1}{Figuras literarias}{Juego de emparejar la figura con su efecto.}

\begin{reto}
«Sus ojos eran dos luceros.» Esta frase es
\opcionesd{una comparación}{una metáfora}{una hipérbole}{una personificación}
\end{reto}

\tecnica{La técnica: ¿qué se dice y qué se quiere decir?}
Una figura retórica dice algo de manera no literal para lograr un efecto (intensificar, embellecer, hacer visible una
idea). Las que más se preguntan:

\begin{center}\small
\begin{tabularx}{\linewidth}{@{}l>{\raggedright\arraybackslash}X>{\raggedright\arraybackslash}X@{}}
\toprule
\textbf{Figura} & \textbf{Qué hace} & \textbf{Ejemplo}\\\midrule
Comparación (símil) & relaciona dos cosas con «como», «cual», «parece» & «blanco como la nieve»\\
Metáfora & identifica una cosa con otra, sin nexo & «tus ojos son luceros»\\
Personificación & da cualidades humanas a lo que no es humano & «el viento susurraba»\\
Hipérbole & exagera & «te lo dije un millón de veces»\\
Antítesis & opone dos ideas & «lo que el día construye, la noche lo deshace»\\
Metonimia & nombra algo por otra cosa relacionada & «leer a García Márquez» (sus obras)\\
Anáfora & repite una palabra al inicio de frases & «Sin ti, sin luz, sin nada»\\\bottomrule
\end{tabularx}
\end{center}

\begin{regla}
\begin{itemize}
  \item «Como» o «parece»: comparación. Sin nexo: metáfora.
  \item Lo no humano que actúa como humano: personificación.
  \item Siempre pregunta qué \emph{significa} la figura en el texto: es lo que más se evalúa.
\end{itemize}
\end{regla}

\begin{texto}{Ciudad de madrugada}{Prosa poética · texto original escrito para este libro}
\pa A las cinco, la ciudad todavía \emph{bosteza}. Los semáforos parpadean para nadie, como si tuvieran sueño, y las
calles son ríos de asfalto que esperan su corriente de carros.

\pa Pasa el primer bus, viejo y ruidoso, \emph{tosiendo} humo en cada esquina. Adentro, las mismas caras de siempre: la
enfermera que vuelve del turno, el panadero que va, el estudiante con un examen en la mochila y un millón de preguntas en
la cabeza.

\pa En una hora, el ruido lo cubrirá todo. Pero ahora, por un momento, la ciudad es un animal grande que respira despacio,
y uno puede escucharle el corazón.
\end{texto}

\subsection{Preguntas comentadas}
\begin{comentada}{«Las calles son ríos de asfalto que esperan su corriente de carros» \pf{1} es}
\mal{A}{una comparación.}{No hay «como» ni otro nexo.}
\bien{B}{una metáfora.}{Identifica las calles con ríos (y los carros con el agua) sin nexo.}
\mal{C}{una hipérbole.}{No exagera una cantidad.}
\mal{D}{una anáfora.}{No hay repetición.}
\end{comentada}

\taller
\piso{Como en el examen}{opción múltiple}
\begin{ejercicios}
\item «La ciudad todavía bosteza» \pf{1} es una
\opcionesd{personificación}{comparación}{metonimia}{antítesis}
\sol{A.}
\item «Un millón de preguntas en la cabeza» \pf{2} es una hipérbole que expresa
\opcionesl{que el estudiante sabe mucho.}{la gran cantidad de dudas o nervios del estudiante.}{que el examen tiene un millón de preguntas.}{que el estudiante está enfermo.}
\sol{B.}
\item «Los semáforos parpadean para nadie, como si tuvieran sueño» combina
\opcionesl{personificación y comparación.}{metáfora y anáfora.}{hipérbole y metonimia.}{solo una antítesis.}
\sol{A.}
\item La imagen final de la ciudad como «un animal grande que respira despacio» transmite
\opcionesl{peligro.}{calma y vida silenciosa antes del ruido del día.}{suciedad.}{tristeza por la ciudad.}
\sol{B.}
\end{ejercicios}

\begin{retoresuelto}
Se identifican los ojos con luceros sin «como»: metáfora. Con «como dos luceros» sería comparación. \textbf{Opción B.}
\end{retoresuelto}
\bitacora

\leccion{pu-lf-2}{Ironía}{Juego de detectar cuándo se dice lo contrario de lo que se piensa.}

\begin{reto}
Después de esperar dos horas el bus bajo la lluvia, Carla dice: «¡Qué servicio tan puntual!». Carla quiere decir que el
servicio
\opcionesd{es excelente}{es impuntual}{es barato}{es nuevo}
\end{reto}

\tecnica{La técnica: el choque entre lo dicho y la situación}
La \clave{ironía} consiste en decir lo contrario de lo que se piensa, o algo distinto, esperando que el lector lo note.
Se detecta por un \emph{choque}: lo que se dice no encaja con la situación (elogiar la puntualidad después de dos horas de
espera). Otras pistas: elogios exagerados, comillas en palabras «positivas», un tono que no corresponde al hecho.

El \clave{sarcasmo} es una ironía hiriente; la \clave{ironía de situación} es cuando los hechos resultan contrarios a lo
esperado (el bombero cuya casa se incendia).

\figura{ironia}{La ironía aparece cuando lo dicho choca con lo que pasa.}

\begin{regla}
\begin{itemize}
  \item Compara lo que se dice con lo que pasa: si chocan, hay ironía.
  \item La intención de la ironía suele ser criticar o burlarse.
  \item En la respuesta, elige lo que el autor \emph{quiere decir}, no lo que dice literalmente.
\end{itemize}
\end{regla}

\begin{texto}{Instrucciones para hacer una fila}{Texto humorístico · texto original escrito para este libro}
\pa Hacer una fila en una oficina pública es un arte que exige preparación. Lo primero es llegar temprano, por ejemplo
dos horas antes de que abran, para descubrir que otras sesenta personas tuvieron la misma brillante idea.

\pa Una vez en la fila, conviene no preguntar nada a nadie, pues cada funcionario tiene la admirable costumbre de enviarlo
a uno a otra ventanilla, donde la fila es todavía más larga y, por supuesto, igual de rápida.

\pa Si al fin llega a la ventanilla, no olvide que le faltará un documento. No importa cuál: siempre falta uno. Esa es la
única certeza del sistema, y hay que agradecerla, porque en un mundo tan cambiante, algo tiene que permanecer.
\end{texto}

\subsection{Preguntas comentadas}
\begin{comentada}{En el \S\,1, la expresión «la misma brillante idea» es irónica porque}
\mal{A}{llegar temprano es realmente la mejor opción.}{Ignora el choque con la situación.}
\bien{B}{la «brillante idea» resulta inútil: sesenta personas llegaron antes.}{El elogio choca con el resultado.}
\mal{C}{el autor admira a los funcionarios.}{\emph{Dice otra cosa}.}
\mal{D}{la fila es corta.}{\emph{Dice lo contrario}.}
\end{comentada}

\taller
\piso{Como en el examen}{opción múltiple}
\begin{ejercicios}
\item «La admirable costumbre» de enviar a otra ventanilla \pf{2} es, en realidad, una crítica a
\opcionesl{la amabilidad de los funcionarios.}{la falta de orientación y la burocracia.}{la rapidez del servicio.}{los ciudadanos.}
\sol{B.}
\item «Igual de rápida» \pf{2} significa, por el contexto, que la otra fila es
\opcionesd{muy rápida}{igual de lenta}{más corta}{inexistente}
\sol{B.}
\item La intención general del texto es
\opcionesl{dar instrucciones reales para hacer filas.}{criticar con humor la ineficiencia de los trámites.}{elogiar el sistema.}{informar sobre horarios.}
\sol{B.}
\end{ejercicios}

\begin{retoresuelto}
«Puntual» choca con dos horas de espera: Carla dice lo contrario de lo que piensa. \textbf{Opción B.}
\end{retoresuelto}
\bitacora

\evaluacion
\begin{texto}{Oda al paraguas olvidado}{Prosa humorística · texto original escrito para este libro}
\pa Nadie ha sido tan fiel como el paraguas que dejé en el bus. Me acompañó tres inviernos sin quejarse, abriéndose
valiente contra aguaceros que habrían hecho llorar a una sombrilla de playa.

\pa Y yo, en agradecimiento, lo abandoné en el asiento de atrás, justo el único día del año en que salió el sol. Qué
manera tan elegante de pagar la lealtad.

\pa Ahora debe viajar por la ciudad, de ruta en ruta, como un huérfano. Espero que alguien lo adopte y lo trate mejor que
yo. Y si un día vuelve a mis manos, le juro que no lo soltaré, aunque llueva.
\end{texto}
\begin{ejercicios}
\item «Abriéndose valiente contra aguaceros» \pf{1} es una
\opcionesd{personificación}{anáfora}{metonimia}{antítesis}
\sol{A. (5.1)}
\item «Qué manera tan elegante de pagar la lealtad» \pf{2} es irónica porque
\opcionesl{el narrador fue realmente elegante.}{abandonar el paraguas es lo contrario de pagar la lealtad.}{el paraguas era elegante.}{el bus era elegante.}
\sol{B. (5.2)}
\item «Como un huérfano» \pf{3} es una
\opcionesd{comparación}{metáfora}{hipérbole}{anáfora}
\sol{A. (5.1)}
\item La frase final, «aunque llueva», produce humor porque
\opcionesl{invierte lo esperado: se usa un paraguas precisamente cuando llueve.}{el narrador odia la lluvia.}{es una hipérbole.}{no tiene sentido.}
\sol{A. Ironía de situación en la promesa. (5.2)}
\end{ejercicios}
